% Options for packages loaded elsewhere
\PassOptionsToPackage{unicode}{hyperref}
\PassOptionsToPackage{hyphens}{url}
\documentclass[
  12pt,
]{article}
\usepackage{xcolor}
\usepackage{amsmath,amssymb}
\setcounter{secnumdepth}{-\maxdimen} % remove section numbering
\usepackage{iftex}
\ifPDFTeX
  \usepackage[T1]{fontenc}
  \usepackage[utf8]{inputenc}
  \usepackage{textcomp} % provide euro and other symbols
\else % if luatex or xetex
  \usepackage{unicode-math} % this also loads fontspec
  \defaultfontfeatures{Scale=MatchLowercase}
  \defaultfontfeatures[\rmfamily]{Ligatures=TeX,Scale=1}
\fi
\usepackage{lmodern}
\ifPDFTeX\else
  % xetex/luatex font selection
\fi
% Use upquote if available, for straight quotes in verbatim environments
\IfFileExists{upquote.sty}{\usepackage{upquote}}{}
\IfFileExists{microtype.sty}{% use microtype if available
  \usepackage[]{microtype}
  \UseMicrotypeSet[protrusion]{basicmath} % disable protrusion for tt fonts
}{}
\makeatletter
\@ifundefined{KOMAClassName}{% if non-KOMA class
  \IfFileExists{parskip.sty}{%
    \usepackage{parskip}
  }{% else
    \setlength{\parindent}{0pt}
    \setlength{\parskip}{6pt plus 2pt minus 1pt}}
}{% if KOMA class
  \KOMAoptions{parskip=half}}
\makeatother
\usepackage{longtable,booktabs,array}
\usepackage{calc} % for calculating minipage widths
% Correct order of tables after \paragraph or \subparagraph
\usepackage{etoolbox}
\makeatletter
\patchcmd\longtable{\par}{\if@noskipsec\mbox{}\fi\par}{}{}
\makeatother
% Allow footnotes in longtable head/foot
\IfFileExists{footnotehyper.sty}{\usepackage{footnotehyper}}{\usepackage{footnote}}
\makesavenoteenv{longtable}
\setlength{\emergencystretch}{3em} % prevent overfull lines
\providecommand{\tightlist}{%
  \setlength{\itemsep}{0pt}\setlength{\parskip}{0pt}}
\usepackage{bookmark}
\IfFileExists{xurl.sty}{\usepackage{xurl}}{} % add URL line breaks if available
\urlstyle{same}
\hypersetup{
  hidelinks,
  pdfcreator={LaTeX via pandoc}}

\author{SCX}
\date{2026}

\begin{document}

\section{Hostile Review Report: Line-by-Line Attack on Situs Theory}\label{hostile-review-report-situs-theory}

\begin{center}\rule{0.5\linewidth}{0.5pt}\end{center}

\textbf{Reviewer Identity}: The most demanding reviewer at Annals of Statistics\\
\textbf{Review Date}: 2026-06-29\\
\textbf{Review Target}: \texttt{paper/situs\_theory/main.tex} + \texttt{theory/self\_evolution/ppe\_rigorous\_derivation.md}\\
\textbf{Review Standard}: Impolite. Uncompromising. Not a single loophole will be spared.

\begin{center}\rule{0.5\linewidth}{0.5pt}\end{center}

\subsection{0. Overall Impression}\label{overall-impression}

This paper claims to have established a ``rigorous mathematical foundation'' for Situs (physical positional encoding) on a framework called SCX. After reading three documents, the conclusion is: \textbf{This paper has fundamental errors or irreparable logical leaps in the proofs of multiple core theorems. The proof of Theorem 2.4.1 is wrong --- not incomplete, wrong. Theorem 1.3.1's Lipschitz constant is off by a factor of approximately 3.5. Theorem 2.2.1's Step 4 is a chasm that cannot be bridged with current tools. Theorem 3$'$'s ``constructive example'' actually falsifies the proposition it claims to prove.} Below is a point-by-point dissection.

\begin{center}\rule{0.5\linewidth}{0.5pt}\end{center}

\subsection{1. Notation Consistency: main.tex vs ppe\_rigorous\_derivation.md}\label{notation-consistency}

\subsubsection{1.1 Inconsistent Theorem Numbering}\label{inconsistent-theorem-numbering}

\begin{longtable}[]{@{}
  >{\raggedright\arraybackslash}p{(\linewidth - 6\tabcolsep) * \real{0.1071}}
  >{\raggedright\arraybackslash}p{(\linewidth - 6\tabcolsep) * \real{0.2321}}
  >{\raggedright\arraybackslash}p{(\linewidth - 6\tabcolsep) * \real{0.5536}}
  >{\raggedright\arraybackslash}p{(\linewidth - 6\tabcolsep) * \real{0.1071}}@{}}
\toprule\noalign{}
\begin{minipage}[b]{\linewidth}\raggedright
Location
\end{minipage} & \begin{minipage}[b]{\linewidth}\raggedright
main.tex number
\end{minipage} & \begin{minipage}[b]{\linewidth}\raggedright
ppe\_rigorous\_derivation.md number
\end{minipage} & \begin{minipage}[b]{\linewidth}\raggedright
Difference
\end{minipage} \\
\midrule\noalign{}
\endhead
\bottomrule\noalign{}
\endlastfoot
Learned PE breakage theorem & Theorem 4.2 (line 1048) & Theorem 4.1 (\S4.2.3, line 851) & \textbf{Off by 1} \\
Encoding imperfection upper bound & Proposition 3.1 (line 864) & Proposition 3.1 (\S3.1.3, line 690) & Consistent \\
Theorem 2$'$ & Theorem 2$'$ (line 881) & Theorem 2$'$ (\S3.2, line 702) & Consistent \\
\end{longtable}

\textbf{Attack}: The numbering gap between Theorem 4.1 and Theorem 4.2 implies a ``vanished Theorem 4.1'' in the paper body. If Theorem 4.1 is Proposition 4.1 (fixed PE preserves Theorem 3), then why isn't Proposition 4.1 called Theorem 4.1? This isn't a typo --- it's a break in the logical chain: Proposition 4.1 uses a trivial deterministic function argument and doesn't deserve the title ``theorem''; but Theorem 4.2's ``destructive'' result is a genuine theorem-level contribution. The numbering chaos exposes the author's own confusion about the importance of their contributions.

\subsubsection{1.2 Split Personality of the Naming System}\label{split-personality-naming}

The paper body uses \texttt{{\textbackslash}Situs}, \texttt{{\textbackslash}Sstates}, \texttt{{\textbackslash}Ppos}, \texttt{{\textbackslash}PE}, \texttt{d\_\{{\textbackslash}pe\}}, while ppe\_rigorous\_derivation.md uses \texttt{PPE}, \texttt{{\textbackslash}mathcal\{S\}}, \texttt{{\textbackslash}mathcal\{P\}}, \texttt{{\textbackslash}text\{PE\}}, \texttt{d\_\{{\textbackslash}text\{pe\}\}}. The paper claims these are ``consistent with'' ppe\_rigorous\_derivation.md (line 1343), yet they differ in every symbol. This is not a LaTeX vs Markdown formatting difference --- the paper uses \texttt{{\textbackslash}Ppos} for position space, while the .md uses \texttt{{\textbackslash}mathcal\{P\}}. The author claims consistency and demonstrates inconsistency in the same sentence.

\subsubsection{1.3 Two Versions of Theorem 2.3.1}\label{two-versions-thm-2.3.1}

\textbf{Paper version} (line 697-701):
\[\delta_s^{\text{PE}} \leq \min\!\left(1 - p_{\text{clean}, s},\, \sqrt{\frac{1}{2} D_{\text{KL}}(P_{\text{PE}|S,Y} \| P_{\text{PE}|S})}\right) + \delta_s^{\text{variance}}\]

\textbf{md version} (line 512-513):
\[\delta_s^{\text{PE}} \leq \min\left(1, \sqrt{\frac{1}{2} \cdot I(P; \text{PE}(P) \mid S, Y)}\right) + \text{Bernstein correction term}\]

These are \textbf{not} the same upper bound. The first uses $1-p_{\text{clean},s}$ and KL divergence (pointwise conditional), the second uses $1$ and mutual information (expectation form). The md version uses ``Bernstein correction'' (implying $O(1/M)$ convergence), the paper uses $\delta_s^{\text{variance}} = O(1/\sqrt{M})$ (implying CLT convergence). Their scaling laws at finite $M$ are completely different --- one is $1/M$, the other is $1/\sqrt{M}$.

\begin{center}\rule{0.5\linewidth}{0.5pt}\end{center}

\subsection{2. Hypothesis Check: Do A1-A6 Still Hold Under Situs?}\label{hypothesis-check-a1-a6-situs}

The CC audit report implicitly assumes 6 core hypotheses for the original SCX theorems. Let us check their status item by item after adding Situs:

\subsubsection{A1: Expert Independence ($v_m \perp v_{m'} \mid s$)}\label{a1-expert-independence}

\textbf{Original status}: $M$ experts' votes on the same state atom $s$ are i.i.d.\ Bernoulli.

\textbf{After adding Situs}: All experts receive the same augmented input $h_i = \phi(s_i) + \text{PE}(p_i)$. The PE term is a shared deterministic signal --- if PE carries information about the label, then all experts' inputs contain the same ``hint.'' This \textbf{weakens} the conditional independence assumption of experts: $\text{Cov}(v_m, v_{m'} \mid s, p) \neq 0$. Strictly speaking, the Chernoff-Hoeffding bound requires conditional independence given the state; now PE injects a common signal shared by all experts. \textbf{This breaks the strictness of the i.i.d. assumption.}

Attack severity: \textbf{Medium}. At large $M$, the correlation introduced by this shared signal may loosen the Chernoff bound by a constant factor (analogous to $\beta$-mixing correction), but the paper does not mention this problem --- it pretends the Chernoff-Hoeffding structure is ``unchanged.''

\subsubsection{A2: Positive Detection Margin ($\Delta_s > 0$)}\label{a2-positive-detection-margin}

\textbf{Original status}: $\Delta_s = p_{\text{noisy},s} - p_{\text{clean},s} > 0$.

\textbf{After adding Situs}: Requires $\Delta_s^{\text{Situs}} = \Delta_s + \delta_s^{\text{PE}} > 0$. But $\delta_s^{\text{PE}} \in [-\Delta_s, 1-p_{\text{clean},s}]$. \textbf{Situs can reduce the detection margin!} If $\delta_s^{\text{PE}} < 0$ and $|\delta_s^{\text{PE}}| > \Delta_s$, then $\Delta_s^{\text{Situs}} < 0$ --- Situs reverses the detection direction. Theorem 2.2.1 only gives sufficient conditions for the existence of some $s$ and some PE making $\delta_s^{\text{PE}} > 0$, but \textbf{does not rule out the possibility of $\delta_s^{\text{PE}} < 0$ for other $s$}.

Attack severity: \textbf{High}. The paper's corrected Theorem 1 replaces $\Delta_s$ with $\Delta_s + \delta_s^{\text{PE}}$ for all states, but $\delta_s^{\text{PE}}$ can be negative. For some $s$, $\exp(-2M(\Delta_s + \delta_s^{\text{PE}})^2)$ may be larger than $\exp(-2M\Delta_s^2)$, i.e., the $F_1$ lower bound \textbf{gets worse}. This is not the ``systematic enhancement'' the paper claims.

\subsubsection{A3: Discreteness of State Space}\label{a3-discreteness-state-space}

\textbf{Original status}: $\mathcal{S} = \{s_1, \ldots, s_N\}$ is a finite discrete set.

\textbf{After adding Situs}: The encoding space $\mathbb{R}^d$ is continuous, and $h_i = \phi(s_i) + \text{PE}(p_i)$ varies continuously in $\mathbb{R}^d$. Even if $s_i$ is discrete, $h_i$ is no longer. This means ``equivalence classes'' among state atoms are no longer strict --- two atoms originally mapping to the same state $s$, if their physical positions differ, will get different $h_i$. This blurs the definition of the probability $\rho_s$ of ``aggregating by state $s$'' --- originally aggregating by state atom type, now needing to aggregate jointly by $(s, p)$.

Attack severity: \textbf{Low-Medium}. Technically fixable (define the distribution on the joint space), but the paper does not discuss this subtlety.

\subsubsection{A4: Bounded Random Variables (Chernoff-Hoeffding Applicability)}\label{a4-bounded-random-variables}

\textbf{Original status}: $v_m \in \{0,1\}$ bounded.

\textbf{After adding Situs}: Votes remain in $\{0,1\}$, Chernoff-Hoeffding applicability is unchanged. This hypothesis \textbf{formally holds}. But as noted in A1, the independence condition is weakened.

Attack severity: \textbf{Low} (formally preserved, substantively implicated by A1).

\subsubsection{A5-A6: Bayes Error Rate and Fano Condition}\label{a5-a6-bayes-fano}

These involve the internal construction of original Theorems 2 and 3, not explicitly numbered in the CC audit report. Under Situs:
- The definition of Bayes error rate in the augmented feature space requires extra care (the augmented space introduces continuous dimensions).
- The $\log|\mathcal{Y}|$ term in Fano's inequality is unchanged, but the computation of conditional entropy involves a continuous-discrete mixture.

Attack severity: \textbf{Low}. These are standard information-theoretic technical details that can be handled, but the paper does not handle them.

\subsubsection{Hypothesis Check Conclusion}\label{hypothesis-check-conclusion}

\textbf{Of the paper's claimed 6 hypotheses, at least A1 (independence) and A2 (positive margin) are significantly affected under Situs. The paper does not verify item by item that these hypotheses still hold before claiming the Chernoff structure is ``unchanged'' --- this is dishonest.}

\begin{center}\rule{0.5\linewidth}{0.5pt}\end{center}

\subsection{3. Proof Completeness: Logical Leaps, Undeclared Assumptions, Erroneous Reasoning}\label{proof-completeness}

\subsubsection{3.1 Theorem 2.2.1 Step 4: Bayes Optimal $\to$ Real Experts Generalization (Fatal)}\label{thm-2.2.1-step4}

Paper body lines 659--661:
\begin{quote}
Step 4 (from Bayes optimal to real experts): Assuming real experts consistently approximate the Bayes optimal (extracting information from augmented features in the same direction as the Bayes optimal), the sign of the above inequality is preserved in expectation.
\end{quote}

The paper itself marks this as \texttt{[Heuristic]}. But this is not heuristic --- \textbf{this is an unbridgeable chasm}. The problems are as follows:

\textbf{(a)} There is no finite-sample convergence guarantee between the behavior of the Bayes optimal classifier given $(s, \text{PE}(p))$ and the behavior of an arbitrary real trained neural network expert $E_m$.

\textbf{(b)} What does ``consistently approximate the Bayes optimal'' mean? If experts are neural networks trained by stochastic gradient descent, their convergence points are some stationary point of the loss landscape, not necessarily the Bayes optimal. Even if they are, it is in the \textbf{population distribution} sense, not guaranteeing the correct sign of $\delta_s^{\text{PE}}$ for every specific state $s$.

\textbf{(c)} The paper claims $p_{\text{clean},s}^{\text{Situs}} \leq p_{\text{clean},s}$ (disagreement on clean samples does not increase). The Bayes optimal classifier indeed does not get worse with more information (guaranteed by the data processing inequality). But real experts may perform \textbf{worse} on clean samples because PE introduces \textbf{spurious correlations} --- especially when training data is limited, PE may learn a pattern that happens to work on the training set but fails on the test set. In this case $p_{\text{clean},s}^{\text{Situs}} > p_{\text{clean},s}$, and $\delta_s^{\text{PE}}$ may be negative.

\textbf{(d)} The entire conclusion of Theorem 2.2.1 --- ``there exists PE and dimension such that $\delta_s^{\text{PE}} > 0$'' --- without Step 4, only holds for the Bayes optimal classifier. For the neural network experts in the actual SCX framework, this is an \textbf{unproven conjecture}.

\textbf{Verdict}: Theorem 2.2.1 is the paper's most important information-theoretic result (connecting $I(Y;P|S) > 0$ with detection improvement). Without Step 4, it only proves that the Bayes optimal classifier can benefit from position information --- which is trivial. Applying this result to real experts in the SCX framework is a \textbf{logical leap}.

\subsubsection{3.2 Theorem 2.4.1 Step 2: Directionality of Fano's Inequality (Fatal)}\label{thm-2.4.1-step2}

This is the most serious error I found. The paper claims Fano gives a \textbf{lower bound} on $\delta_s^{\text{PE}}$. But examining the derivation closely:

\textbf{Step 1}: Fano's inequality standard form (correct):
\[H(Y \mid S, \text{PE}(P)) \leq H_2(P_e) + P_e \cdot \log(|\mathcal{Y}| - 1)\]

From this (paper lines 750--752):
\[P_e \geq \frac{H(Y \mid S, \text{PE}(P)) - 1}{\log |\mathcal{Y}|}\]

This is a correct \textbf{lower bound}.

\textbf{Step 3} (paper lines 763--766, md lines 601--607):
\[P_e - P_e^{\text{Situs}} = \frac{H(Y|S) - H(Y|S, \text{PE}(P))}{\log |\mathcal{Y}|} \quad (\text{taking equality in Fano})\]

\textbf{This is a fundamental error.} The logical chain is:
1. $P_e \geq \frac{H(Y|S)-1}{\log|\mathcal{Y}|}$ (Fano)
2. $P_e^{\text{Situs}} \geq \frac{H(Y|S,\text{PE}(P))-1}{\log|\mathcal{Y}|}$ (Fano)
3. If both inequalities attain equality $\to$ $P_e - P_e^{\text{Situs}} = \frac{H(Y|S) - H(Y|S,\text{PE}(P))}{\log|\mathcal{Y}|}$

But to go from the \textbf{inequalities} (1) and (2) to the \textbf{equality} (3), one needs (i) Fano's inequality to attain equality in both worlds, and (ii) ``the difference of two Fano equalities = the difference of error rates.''

The problem:
- Fano's inequality only attains equality for specific distributions (noise uniformly distributed across error classes, and $H_2(P_e)$ being binary entropy), which does not hold in general.
- More importantly, even if both Fano bounds attain equality, $a \geq A$ and $b \geq B$ \textbf{does not imply} $a-b = A-B$.

\textbf{The essence of the directionality error}: The paper attempts to derive a \textbf{lower bound} on error rate improvement from Fano's \textbf{upper bound} (conditional entropy). Fano's inequality has the form $H(Y|X) \leq f(P_e)$, giving a lower bound on $P_e$. But $P_e - P_e^{\text{Situs}}$ as a difference cannot be directly derived from two lower bounds on $P_e$ --- the direction of the difference is not controlled by individual lower bounds.

\textbf{Correct approach}: One needs an \textbf{upper bound} on $P_e - P_e^{\text{Situs}}$ to use Fano for any rigorous bound. Or use stronger tools such as the inverse form of Fano (but this typically requires additional distributional assumptions).

\textbf{Verdict}: \textbf{The proof of Theorem 2.4.1 is wrong.} Not incomplete, not heuristic --- it is a logical fallacy of deducing an equality from inequalities. The word ``lower bound'' in the theorem statement misleads readers into thinking this is a rigorous information-theoretic bound. At best, it is a heuristic guess.

\subsubsection{3.3 Theorem 2.2.1 Step 3: ``Clean Sample Disagreement Does Not Increase'' Argument}\label{thm-2.2.1-step3}

Paper line 655:
\begin{quote}
For clean samples, additional position information further confirms the correct label $\to$ disagreement probability does not increase: $p_{\text{clean}, s}^{\text{Situs}} \leq p_{\text{clean}, s}$.
\end{quote}

This relies on an implicit assumption: PE's information is ``consistent with'' rather than ``conflicting with'' the correct label. But:
- PE is derived from physical position, which may have \textbf{spurious correlations} with the label (correlated in the training distribution but not in the target distribution).
- PE may amplify some systematic bias of experts --- if all experts share a common blind spot, additional PE information may make them collectively more confidently wrong on that blind spot.
- The paper itself acknowledges this risk: in \S6.3's ``one-vote veto condition,'' F5 points out that under ``random or pseudo-random configurations + finite samples,'' $\delta_s^{\text{PE}} > 0$ on the training set may turn into $\delta_s^{\text{PE}} < 0$ on the test set.

\textbf{Attack}: The Step 3 argument assumes PE's information is always ``beneficial'' --- it increases disagreement on noisy samples but not on clean samples. This is equivalent to \textbf{assuming the conclusion}: PE somehow distinguishes noisy and clean samples. If PE's information is a purely additional feature (correlational but not causal), it may increase disagreement on both types of samples by the same magnitude, leading to $\delta_s^{\text{PE}} = 0$.

\subsubsection{3.4 Theorem 2.5.1: Illegitimate Use of Lipschitz Continuity Assumption}\label{thm-2.5.1-lipschitz}

Theorem 2.5.1 claims:
\[|\delta_i^{\text{PE}} - \delta_j^{\text{PE}}| \leq 2 \cdot L_E \cdot L_{\text{PE}} \cdot d_{\mathcal{P}}(p_i, p_j)\]

The proof uses (md line 640):
\[|p_{\text{clean}, i}^{\text{PPE}} - p_{\text{clean}, j}^{\text{PPE}}| \leq L_E \cdot \|h_i - h_j\|\]

But $p_{\text{clean}, i}^{\text{PPE}} = \mathbb{E}[\mathbf{1}[E_m(h_i) \neq y]]$ is the \textbf{expectation of an indicator function}. Even if $E_m$'s output is $L_E$-Lipschitz (in some sense), the indicator function $\mathbf{1}[\cdot \neq y]$ is \textbf{discontinuous} --- it jumps at the decision boundary. The $L_E$-Lipschitz assumption typically applies to the network's output logits or softmax probabilities, not the expectation of a 0/1 decision.

For this bound to hold, additional assumptions are needed:
- The expert's decision boundary is sufficiently far from both $h_i$ and $h_j$ (margin condition), or
- The expert outputs soft probabilities and the expectation of the indicator function happens to be Lipschitz (requiring the expert's predicted probabilities to be Lipschitz in the input, and the threshold-crossing probability to be a zero-measure set).

The paper provides none of these conditions. \textbf{This is not a technical detail --- it is the key to whether the bound is non-vacuous.}

\begin{center}\rule{0.5\linewidth}{0.5pt}\end{center}

\subsection{4. Notation Errors: Not Just One}\label{notation-errors}

\subsubsection{4.1 Sign Error in CC Audit Report (Corrected, Confirmed)}\label{sign-error-cc}

The sign in Proposition 2.1 of the CC audit report was indeed wrong. Its $\delta_s^{\text{PE}} = (p_{\text{clean}}^{\text{PPE}} - p_{\text{clean}}) - (p_{\text{noisy}}^{\text{PPE}} - p_{\text{noisy}})$ has the physical meaning: encoding is helpful only when the consistency gain on clean samples exceeds the ``confusion gain'' on noisy samples. This is counter to the intuition that ``encoding should increase the noise detection margin.'' The paper's correction (swapping the order of the two terms) is correct.

\textbf{However}: The CC audit report's internal notation is itself inconsistent. \S0.2 defines $p_{\text{clean}}$ as the disagreement probability ($<0.5$), but \S2.3's proof contains formulas like $|p_{\text{clean}} + p_{\text{noisy}} - 1|$ that treat $p_{\text{clean}}$ as a consistency probability. So the CC report's ``sign error'' may not be a simple sign flip --- it is a deeper \textbf{internal schism in the notation system}: \S0.2 and \S2.3 use completely different semantics for $p$.

The paper corrected this sign error but \textbf{did not point out the internal notation schism in the CC report} --- it simply flipped the sign and claimed a correction. In reality, the CC report's problem is more severe.

\subsubsection{4.2 Theorem 1.3.1: Massive Overestimation of Lipschitz Constant (New Derivation Error)}\label{thm-1.3.1-overestimation}

This is not a notation error --- it is an order-of-magnitude error. The paper claims:
\[L_{\text{PE}}^{\text{rot}} = 2\sqrt{\alpha^2 + \beta^2 + \gamma^2}\]

The actual Lipschitz constant is $\max(\alpha, \beta, \gamma)$. Off by a factor of approximately $2\sqrt{3} \approx 3.46$ (when $\alpha=\beta=\gamma$).

\textbf{Derivation}: For $\mathbf{e}_0 = (1,0,0,1,0,0,\ldots)^T$ (first dimension of each rotation plane is 1),
\[\|\text{PE}_{\text{rot}}(\mathbf{p}) - \text{PE}_{\text{rot}}(\mathbf{q})\|^2 = 4\sin^2(\alpha\Delta x/2) + 4\sin^2(\beta\Delta y/2) + 4\sin^2(\gamma\Delta z/2)\]

As $\|\Delta\mathbf{p}\| \to 0$:
\[\|\text{PE}_{\text{rot}}(\mathbf{p}) - \text{PE}_{\text{rot}}(\mathbf{q})\| \approx \sqrt{\alpha^2\Delta x^2 + \beta^2\Delta y^2 + \gamma^2\Delta z^2}\]

By Cauchy-Schwarz:
\[\sqrt{\alpha^2\Delta x^2 + \beta^2\Delta y^2 + \gamma^2\Delta z^2} \leq \max(\alpha, \beta, \gamma) \cdot \|\Delta\mathbf{p}\|_2\]

\textbf{The true Lipschitz constant is $\max(\alpha, \beta, \gamma)$}, and this bound is tight (take $\Delta\mathbf{p}$ along the direction of the maximum frequency parameter).

The error in the paper's proof lies in Step 2's triangle inequality decomposition (md lines 300--317). The author correctly obtains:
\[\|\mathbf{R}(\mathbf{p}) - \mathbf{R}(\mathbf{q})\|_F \leq 2\alpha|\Delta x| + 2\beta|\Delta y| + 2\gamma|\Delta z|\]

But the three difference matrices act on \textbf{mutually disjoint subspaces} (dimension pairs (0,1), (2,3), (4,5)). For a sum of matrices on mutually disjoint subspaces $D = D_x + D_y + D_z$:
\[\|D\|_F^2 = \|D_x\|_F^2 + \|D_y\|_F^2 + \|D_z\|_F^2\]

The triangle inequality $\|D\|_F \leq \|D_x\|_F + \|D_y\|_F + \|D_z\|_F$ can be loose by a factor of $\sqrt{3}$ in the worst case (three equal-norm terms). Then the author introduces another factor of $\sqrt{3}$ looseness through Cauchy-Schwarz ($\ell_1 \to \ell_2$ conversion).

\textbf{The two loosenings multiply to give a final constant roughly 3.46 times too high.} Worse, the author noticed this problem (Remark 1.3.1/md note 1.3.1 acknowledges that the effective constant for a single direction is only $\alpha$, not $2\alpha$), yet still wrote the loose bound as the ``exact constant'' in the theorem statement.

\textbf{Verdict}: The ``exact Lipschitz constant'' claimed in Theorem 1.3.1 is neither exact nor tight. The true exact Lipschitz constant (using $\|\cdot\|_2$ vector norm and $\|\cdot\|_2$ position norm) is $\max(\alpha, \beta, \gamma)$.

\subsubsection{4.3 Missing Normalization Factor in Theorem 1.2.1}\label{missing-normalization-thm-1.2.1}

The encoding kernel definition (Definition 1.2.1) is:
\[\langle \text{PE}_{\text{scalar}}(p), \text{PE}_{\text{scalar}}(q) \rangle = \sum_{j=0}^{d/2-1} \cos(2\pi(p-q)/\lambda_j) = K_{\text{PE}}(\Delta)\]

At $\Delta = 0$, $K_{\text{PE}}(0) = d/2$. But the target kernel $k(0) = 1$.

As $d \to \infty$, the encoding kernel \textbf{diverges} to infinity rather than converging to the target kernel. To make the encoding kernel approximate the target kernel, normalization is required:
\[\frac{2}{d} K_{\text{PE}}(\Delta) \to \int_0^\infty S(\omega)\cos(\omega\Delta)d\omega = k(\Delta)\]

But Theorem 1.2.1's statement and proof \textbf{completely omit} this $2/d$ factor. The theorem's statement that ``the encoding kernel optimally approximates $k(\Delta)$ in $L^2([0,L])$'' is mathematically wrong --- $K_{\text{PE}}$ and $k$ are at different scales ($O(d)$ vs $O(1)$).

\subsubsection{4.4 Theorem 2.3.1: Legitimacy of Pinsker Application}\label{pinsker-legitimacy}

Theorem 2.3.1 uses Pinsker's inequality to relate $\delta_s^{\text{PE}}$ to KL divergence. But there is a subtle issue in the use of Pinsker's inequality:

Pinsker's bound: $|P(A) - Q(A)| \leq \sqrt{\frac{1}{2}D_{\text{KL}}(P\|Q)}$, where $P$ and $Q$ are \textbf{probability measures on the same measurable space} and $A$ is an event in that space.

In the paper's application, $P$ and $Q$ are the conditional distributions of $\text{PE}(P)$ given $S$ (clean vs mixed), but the quantity being bounded, $p_{\text{clean},s}^{\text{Situs}} - p_{\text{clean},s}$, involves \textbf{the expert's output event}, which is not in the space of $\text{PE}(P)$.

Strictly, the correct argument is:
\[p_{\text{clean},s}^{\text{Situs}} = \int P(E_m(\phi(s)+z) \neq y) \cdot dP_{\text{PE}(P)|S=s,\text{clean}}(z)\]
\[p_{\text{clean},s} = P(E_m(\phi(s)) \neq y)\]

The second quantity does not even involve the distribution of $\text{PE}(P)$. The paper compares $p_{\text{clean},s}^{\text{Situs}}$ and $p_{\text{clean},s}$, but Pinsker only constrains changes in the distribution of $\text{PE}(P)$, not directly constrain changes in output probability --- unless one additionally assumes the expert's output dependence on PE is Lipschitz (returning to the Theorem 2.5.1 problem).

\textbf{However, this may be fixable}: If we compare $p_{\text{clean},s}^{\text{Situs}}$ and $p_{\text{clean},s}^{\text{avg}} = \int P(E_m(\phi(s)+z) \neq y) \cdot dP_{\text{PE}(P)|S=s}(z)$ (i.e., the clean disagreement rate averaged over PE), Pinsker can be legitimately applied. But the paper does not make this decomposition.

\begin{center}\rule{0.5\linewidth}{0.5pt}\end{center}

\subsection{5. Tightness: Claimed ``Tight'' but None Are Tight}\label{tightness-none-tight}

\subsubsection{5.1 Theorem 1.2.3: Scalar Sine Lipschitz Constant}\label{thm-1.2.3-tight}

Claimed ``tight'' (line 414, md line 203). The derivative of each $(\sin,\cos)$ pair indeed satisfies $\cos^2 + \sin^2 = 1$ identically for all $p$, so for any two arbitrarily close points, the local rate of change of this pair indeed achieves $\omega_j$. For the entire vector, all pairs act independently in orthogonal subspaces, so the Lipschitz constant is indeed $\sqrt{\sum_j \omega_j^2 \cdot 2}$ = the claimed value. \textbf{This bound is indeed tight.} \checkmark

\subsubsection{5.2 Theorem 1.3.1: 3D Rotation Lipschitz Constant}\label{thm-1.3.1-tight}

Claimed ``worst-case'' tight (line 519, md line 337). As shown in \S4.2, \textbf{not tight} --- off by approximately 3.46$\times$. The true tight Lipschitz constant is $\max(\alpha, \beta, \gamma)$. $\times$

\subsubsection{5.3 Theorem 1.2.2: Approximation Error Bound}\label{thm-1.2.2-tight}

Claims $O(1/d)$ with constant $\frac{2\pi^2\xi}{d} \cdot \frac{L}{\xi} = \frac{2\pi^2 L}{d}$. \textbf{No proof.} Moreover, $\xi$ appears in both numerator and denominator and cancels, indicating that the constant $\frac{2\pi^2 L}{d}$ is independent of the physical correlation length $\xi$ --- this is highly suspicious. How can a kernel's approximation quality be independent of the kernel's width?

Actual derivation (mine): For Riemann-Stieltjes sum approximation of $\cos(\omega\Delta)$ on $[0,L]$ using quantile sampling, the error is jointly controlled by truncated high-frequency components and the coarsest low-frequency resolution. At fixed $d$, the error should vary with $\xi$ --- small $\xi$ (short-range correlation) means more high-frequency content, harder to approximate. The paper's claimed $\xi$-independent constant is physically absurd.

\subsubsection{5.4 Theorem 2$'$'s Bound}\label{thm-2prime-tight}

Claims $F_{1,\text{SCX}}^{\text{Situs}} \leq F_{1,\text{base}} + C_F \cdot \sqrt{(\delta + 2\varepsilon_{\text{PE}}/C_F^2)/2}$.

Note that as $\varepsilon_{\text{PE}} \to I(Y;P|X)$ (completely failed encoding):
\[\delta_{\text{PPE}} = \delta + \frac{2I(Y;P|X)}{C_F^2}\]

The upper bound is looser than the original Theorem 2 by a factor of $\sqrt{\delta + 2I/C_F^2} / \sqrt{\delta}$. When $I(Y;P|X)$ is large (position has lots of information but the encoding completely loses it), the upper bound can become arbitrarily loose. This is certainly not a ``tight'' bound --- it is a \textbf{non-vacuous but potentially very loose} upper bound. The paper does not claim it is tight (honest on this point), but the abstract says ``corrected weak feature failure upper bound'' implying some precision, when in reality it just adds a positive term to an already loose bound.

\begin{center}\rule{0.5\linewidth}{0.5pt}\end{center}

\subsection{6. Theorem 3$'$'s Constructive Example: Shooting Itself in the Foot}\label{thm-3prime-constructive-example}

\subsubsection{6.1 Contradiction Between Claim and Reality}\label{contradiction-claim-reality}

\S4.3 (paper lines 1096--1137) is titled ``Constructive Minimal Example.'' What does this ``constructive minimal example'' demonstrate?

\begin{enumerate}
\def\labelenumi{\arabic{enumi}.}
\tightlist
\item
  Constructs two worlds $W_A, W_B$, verifying $P_{W_A}(X,Y) = P_{W_B}(X,Y)$ \checkmark
\item
  Introduces learned PE $\text{PE}_\theta(p) = \theta \cdot p$
\item
  Computes $\mathbb{E}_{W_A}[Y|X] = \mathbb{E}_{W_B}[Y|X] = 0.6 \cdot \text{sign}(x)$ --- \textbf{numerically identical}
\item
  Conclusion: \textbf{Even learned PE cannot distinguish the two worlds} (paper lines 1131--1132 original text)
\end{enumerate}

\textbf{This is a ``constructive counterexample'' --- it constructs an example where learned PE cannot distinguish, exactly proving Theorem 3's robustness rather than its fragility!}

\subsubsection{6.2 Hand-Waving ``Conditions That Truly Enable Distinction''}\label{hand-waving-conditions}

After proving that learned PE cannot distinguish, the paper (lines 1134--1136) waves its hands and says:
\begin{quote}
``However, when the auxiliary loss involves unobserved physical quantities $Z$ --- for example, in $W_A$ noisy labels cause $|Z_{\text{pred}} - Z_{\text{DFT}}|$ to be large (physical inconsistency), while in $W_B$ this difference is small (difficult but physically consistent) --- then the optimal parameters of learned PE differ between the two worlds, and indistinguishability is broken.''
\end{quote}

This is \textbf{not a constructive proof} --- it is a verbal description of an imagined scenario. No concrete definition of $Z$, no concrete form of $\mathcal{L}_{\text{aux}}$, no verification that the two worlds indeed give different optimal points. The entire \S4.3 promises a ``constructive minimal example'' but delivers only an \textbf{empty claim}.

\subsubsection{6.3 Final Verdict on Theorem 3$'$}\label{final-verdict-thm-3prime}

Theorem 3$'$ itself is a reasonable conditional theorem statement (adding premises (A)(B) preserves indistinguishability), but:
- The ``distinguishability'' conclusion when premise (A) fails \textbf{lacks constructive proof support}.
- The ``constructive minimal example'' is actually a \textbf{counterexample to the counterexample} --- i.e., it demonstrates the theorem's robustness rather than the breakage direction.
- The \S4.3 content contradicts the paper's claimed contribution.

\end{document}
